\emph{Preliminaries.} From the Definition: physical vectors $g_{b} = \hb^{-1}(b)$ and $g^{\mu} = \hb^{T}(e^{\mu})$ (clause (a)); $D_{a}M = a \cdot \nabla M + \Omega(a) \times M$, $A \cdot D = D_{\hb(A)}$ and $D = \hb^{T}(\partial_{a})D_{a}$ (clause (b)); $R(a \wedge b) = a \cdot \nabla\,\Omega(b) - b \cdot \nabla\,\Omega(a) + \Omega(a) \times \Omega(b)$ and $\mathcal{R}(B) = -R(\hb(B))$ (clause (c)); the gauge freedoms (clause (e)); the torsion equation $D \wedge \hb^{T}(a) = \hb^{T}(\nabla \wedge a)$ (clause (f)). Its field equation is not used. For a bivector $\mathcal{B}$ and vectors $x, y$: $\mathcal{B} \times y = \mathcal{B} \cdot y = -y \cdot \mathcal{B}$ and $x \cdot (\mathcal{B} \cdot y) = -y \cdot (\mathcal{B} \cdot x)$. $D_{a}$ is a derivation of every product of multivector fields. \emph{The torsion equation, two forms.} (T1) It holds for every vector field $a$ exactly when $D \wedge \hb^{T}(\nabla\phi) = 0$ for every scalar field $\phi$: one way because $\nabla \wedge \nabla\phi = 0$; the other because $a = \sum_{\mu}a_{\mu}\nabla x^{\mu}$ and $D \wedge (\phi X) = \hb^{T}(\nabla\phi) \wedge X + \phi\,D \wedge X$, so $D \wedge \hb^{T}(a) = \sum_{\mu}\hb^{T}(\nabla a_{\mu}) \wedge \hb^{T}(e^{\mu}) = \hb^{T}(\nabla \wedge a)$. (T2) For constant $a, b$ it gives $D_{a}g_{b} = D_{b}g_{a}$: at $\phi = x^{\mu}$, $g^{\nu} \wedge D_{e_{\nu}}g^{\mu} = 0$; contracting with $g_{a} \wedge g_{b}$ and using $g_{a} \cdot g^{\nu} = a \cdot e^{\nu}$ gives $g_{a} \cdot D_{b}g^{\mu} = g_{b} \cdot D_{a}g^{\mu}$, and since $g_{a} \cdot g^{\mu}$ is constant, $(D_{b}g_{a}) \cdot g^{\mu} = -g_{a} \cdot D_{b}g^{\mu}$; so $(D_{a}g_{b} - D_{b}g_{a}) \cdot g^{\mu} = 0$ for every $\mu$. If $\Omega(a) \cdot b = \Omega(b) \cdot a$ for all $a, b$ then $\Omega = 0$ (the form $c \cdot (\Omega(a) \cdot b)$ is symmetric in $a, b$ and antisymmetric in $b, c$), so (T2) fixes $\Omega$ from $\hb$ wherever $\hb$ is invertible, as a rational function of $\hb^{-1}$ and its first derivatives. \emph{The family.} On a bounded region $\hb_{\varepsilon} = \id + \varepsilon H_{1} + \varepsilon^{2}H_{2}$ is invertible for small $\varepsilon$, and $\Omega_{\varepsilon}$, $\mathcal{R}_{\varepsilon}$, $\mathcal{G}_{\varepsilon}$ are smooth in $\varepsilon$ there, with $\Omega_{0} = 0$ and $\mathcal{G}_{0} = 0$. An identity that holds for every field obeying the torsion equation holds for every small $\varepsilon$, hence for each coefficient of its expansion. The coefficient of $\varepsilon^{2}$ of $\mathcal{G}_{\varepsilon}$ is linear in $H_{2}$, through the same linear map as the coefficient of $\varepsilon$ in $H_{1}$ (replace $\varepsilon$ by its square and $H_{1}$ by $H_{2}$), plus a part quadratic in $H_{1}$; $\Omega$ has one derivative of $\hb$ in each term and $R$ one derivative of $\Omega$ or two factors $\Omega$, so each term of $\mathcal{G}^{(2)}[H_{1}]$ has two derivatives.

\emph{(i).} With $g_{b} = b - \varepsilon H_{1}(b) + O(\varepsilon^{2})$, the coefficient of $\varepsilon$ of (T2) is $\Omega_{1}(a) \cdot b - \Omega_{1}(b) \cdot a = (a \cdot \nabla)H_{1}(b) - (b \cdot \nabla)H_{1}(a)$, which fixes $\Omega_{1}$. The stated $\Omega_{1}$ satisfies it: $(\dot\nabla \wedge \dot S(a)) \cdot b = \dot\nabla(\dot S(a) \cdot b) - (b \cdot \nabla)S(a)$, whose antisymmetric part in $a, b$ is $(a \cdot \nabla)S(b) - (b \cdot \nabla)S(a)$ because $S$ is symmetric; and $-((a \cdot \nabla)A) \cdot b + ((b \cdot \nabla)A) \cdot a = (a \cdot \nabla)(b \cdot A) - (b \cdot \nabla)(a \cdot A)$. (That $H_{1}(a) = S(a) + a \cdot A$: $\partial_{a} \wedge (a \cdot A) = 2A$ and $\partial_{a} \wedge S(a) = 0$.) At first order $R(a \wedge b) = (a \cdot \nabla)\Omega_{1}(b) - (b \cdot \nabla)\Omega_{1}(a)$, in which the terms with $A$ cancel, second derivatives commuting; what is left is the field strength of $\Omega(a) = -\nabla \wedge s(a)$ with $s = -S$, the weak symmetric field of the proof of clause (vii) of Theorem \claimref{synthesis.field-equation-linear-in-the-curvature-with-a-conserved-source-is-einsteins}, whose curvature operator is the stored $\mathcal{R}[h]$ with $h = 2s = -(H_{1} + H_{1}^{T})$ and whose Einstein map is $\mathcal{E}[h]$; $\hb$ is the identity at this order on the argument of $\mathcal{R}$. The derived metric is $(a - \varepsilon H_{1}(a)) \cdot (b - \varepsilon H_{1}(b)) = a \cdot b + \varepsilon\,h(a) \cdot b$ to first order. The formula for $\mathcal{E}$ is the contraction of the stored curvature, as in the proof of Theorem \claimref{synthesis.null-plane-data-of-the-first-order-gravitational-field-are-the-shear-of-the-generators}; each of its groups is a symmetric map, and its flat divergence vanishes by clause (v) of that theorem (the values of the stored Einstein map are conserved) with that symmetry. $\underline{\xi} + \underline{\xi}^{T}$ is the displacement $a \mapsto (a \cdot \nabla)\xi + \nabla(a \cdot \xi)$, on which the stored first-order curvature vanishes (Theorem \claimref{synthesis.weyl-flat-variations-carry-only-the-stress-trace-and-the-whole-stress-clause-is-the-trace-free-channel}(i)), and $\beta_{B} + \beta_{B}^{T} = 0$. The check computes the coefficient of $\varepsilon$ from the gauge fields for $H_{1} = -\tfrac14(p \cdot x)^{2}h_{0}$, with and without an antisymmetric part, and finds the stored $\mathcal{G}_{p}[h_{0}]$, formed from the stored curvature, in all sixteen components; on a polynomial $H_{1}$ that is not symmetric it finds $\mathcal{E}[-(H_{1} + H_{1}^{T})]$ and the stated $\Omega_{1}$.

\emph{(ii).} By clause (i) of the second cited theorem the curvature operator of every field obeying the torsion equation satisfies the cyclic identity, and by its clause (ii) every such operator has a symmetric Ricci map; the Einstein map is then symmetric. So $\mathcal{G}_{\varepsilon}$ is symmetric for every small $\varepsilon$, and so is its coefficient $\mathcal{G}^{(1)}[H_{2}] + \mathcal{G}^{(2)}[H_{1}]$; the first term is symmetric by (i). Nothing was assumed of $H_{1}$. The same holds for $\mathcal{G}^{(2)}_{r}$. (If the Einstein map is transcribed to the background by maps that depart from the identity at first order, $P_{\varepsilon}\,\mathcal{G}_{\varepsilon}\,Q_{\varepsilon}$, the quadratic part acquires $P_{1}\mathcal{G}^{(1)}[H_{1}] + \mathcal{G}^{(1)}[H_{1}]Q_{1}$, which need not be symmetric and vanishes on first-order solutions.)

\emph{(iii).} By clause (iii) of the second cited theorem $\operatorname{div}\mathcal{G}_{\varepsilon} = 0$ for every small $\varepsilon$, where for a field $M$ of maps of physical vectors and constant vectors $a, c$, $(a \cdot D\,M)(c) = D_{\hb(a)}(M(c)) - M(D_{\hb(a)}c)$ and $(\operatorname{div}M)(c) = \partial_{a} \cdot ((a \cdot D\,M)(c))$. Since $c$ is constant, $D_{\hb(a)}c = \Omega(\hb(a)) \cdot c$, so
\[
(\operatorname{div}M)(c) = \partial_{a} \cdot \Bigl((\hb(a) \cdot \nabla)M(c) + \Omega(\hb(a)) \cdot M(c) - M\bigl(\Omega(\hb(a)) \cdot c\bigr)\Bigr) .
\]
Put $M = \varepsilon M_{1} + \varepsilon^{2}M_{2}$, $\hb = \id + \varepsilon H_{1}$ and $\Omega(\hb(a)) = \varepsilon\,\Omega_{1}(a)$, each up to higher order. The coefficient of $\varepsilon$ is $\nabla \cdot (M_{1}(c))$. The coefficient of $\varepsilon^{2}$ is the sum of four terms. First, $\nabla \cdot (M_{2}(c))$. Second, with $u = M_{1}(c)$, $\partial_{a} \cdot ((H_{1}(a) \cdot \nabla)u) = \sum_{a}(H_{1}(e_{a}) \cdot \nabla)(e^{a} \cdot u) = \dot\nabla \cdot H_{1}(\dot u)$, since $H_{1}(u) = \sum_{a}(e^{a} \cdot u)H_{1}(e_{a})$. Third, $\partial_{a} \cdot (\Omega_{1}(a) \cdot u) = -u \cdot (\Omega_{1}(\partial_{a}) \cdot a) = u \cdot (\partial_{a} \cdot \Omega_{1}(a)) = u \cdot w$. Fourth, for symmetric $M_{1}$, $-\partial_{a} \cdot M_{1}(\Omega_{1}(a) \cdot c) = -M_{1}(\partial_{a}) \cdot (\Omega_{1}(a) \cdot c) = c \cdot (\Omega_{1}(\partial_{a}) \cdot M_{1}(a))$. So
\[
(\operatorname{div}M)^{(2)}(c) = \nabla \cdot (M_{2}(c)) + \dot\nabla \cdot H_{1}(\dot u) + u \cdot w + c \cdot \bigl(\Omega_{1}(\partial_{a}) \cdot M_{1}(a)\bigr) . \tag{$*$}
\]
For $M = \mathcal{G}_{\varepsilon}$: $M_{1} = \mathcal{G}^{(1)}[H_{1}]$, symmetric by (i), and $M_{2} = \mathcal{G}^{(1)}[H_{2}] + \mathcal{G}^{(2)}[H_{1}]$, whose first term has no flat divergence by (i). Setting $(*)$ to zero gives the stated identity; every term on its right has $\mathcal{G}^{(1)}[H_{1}]$ or a first derivative of it as a factor. The check verifies the identity exactly on a polynomial $H_{1}$ that is neither symmetric nor a solution, each of the twelve terms and all four left sides nonzero, and that leaving out any one of the three terms fails; and on two first-order solutions, built from three null polynomial waves with transverse trace-free polarisations, a polynomial displacement and a polynomial rotation gauge, that the whole jet of $\mathcal{G}^{(1)}[H_{1}]$ vanishes and the flat divergence of $\mathcal{G}^{(2)}[H_{1}]$ vanishes in all four directions, as a cancellation of sixteen nonzero terms.

\emph{(iv).} \emph{Covariance.} (1) For a rotor field $U$ let $\hb'(a) = \hb(\tilde U a U)$ and $\Omega'(a) = U\,\Omega(a)\,\tilde U - 2(a \cdot \nabla U)\tilde U$ (clause (e)). Then $\hb'^{T}(a) = U\,\hb^{T}(a)\,\tilde U$, and for a multivector field $K$, $D'_{a}(UK\tilde U) = U(D_{a}K)\tilde U$: indeed $a \cdot \nabla(UK\tilde U) = U(a \cdot \nabla K)\tilde U + (a \cdot \nabla U)K\tilde U + UK(a \cdot \nabla\tilde U)$, while $-2(a \cdot \nabla U)\tilde U \times (UK\tilde U) = -(a \cdot \nabla U)K\tilde U + UK\tilde U(a \cdot \nabla U)\tilde U = -(a \cdot \nabla U)K\tilde U - UK(a \cdot \nabla\tilde U)$, using $\tilde U(a \cdot \nabla U) = -(a \cdot \nabla\tilde U)U$. With $K = \hb^{T}(\nabla\phi)$, $D' \wedge \hb'^{T}(\nabla\phi) = U\,(D \wedge K)\,\tilde U = 0$, so by (T1) the primed pair obeys the torsion equation: $\Omega'$ is the rotation gauge field that the torsion equation assigns to $\hb'$. By clause (e), $\mathcal{R}'(B) = U\,\mathcal{R}(\tilde U B U)\,\tilde U$; contracting, $\mathcal{G}'(a) = U\,\mathcal{G}(\tilde U a U)\,\tilde U$. (2) For a displacement $f$ with differential $\underline{f}$ and its adjoint $\bar f$, let $\hb'(a; x) = \underline{f}^{-1}(\hb(a; f(x)))$ and $\Omega'(a; x) = \Omega(\underline{f}(a); f(x))$ (clause (e)). For a scalar field $\phi$ put $\phi'(x) = \phi(f(x))$; then $\nabla\phi' = \bar f(\nabla\phi)$ at $f(x)$, $\hb'^{T} = \hb^{T}\bar f^{-1}$, so $\hb'^{T}(\nabla\phi')(x) = K(f(x))$ with $K = \hb^{T}(\nabla\phi)$, and $D'_{a}(K \circ f) = (D_{\underline{f}(a)}K) \circ f$ by the chain rule. Hence $D' \wedge \hb'^{T}(\nabla\phi') = \hb^{T}(\bar f^{-1}(\partial_{a})) \wedge (D_{\underline{f}(a)}K) \circ f = (D \wedge K) \circ f = 0$, and every scalar field is such a $\phi'$: the primed pair obeys the torsion equation. For the field strength, $a \cdot \nabla\,[\Omega(\underline{f}(b); f(x))] = (\underline{f}(a) \cdot \nabla\,\Omega)(\underline{f}(b)) + \Omega((a \cdot \nabla)(b \cdot \nabla)f)$, the last term symmetric in $a, b$; so $R'(a \wedge b; x) = R(\underline{f}(a) \wedge \underline{f}(b); f(x))$, $\mathcal{R}'(B; x) = -R'(\hb'(B)) = \mathcal{R}(B; f(x))$ and $\mathcal{G}'(a; x) = \mathcal{G}(a; f(x))$.

\emph{Rotation.} With $U = e^{\varepsilon B/2}$, $\tilde U a U = a + \varepsilon\,\beta(a) + \tfrac12\varepsilon^{2}\beta(\beta(a)) + O(\varepsilon^{3})$, $\beta = \beta_{B}$ (the derivative in $\varepsilon$ of $\tilde U a U$ is $(\tilde U a U) \cdot B$). So $\hb'_{\varepsilon} = \id + \varepsilon(H_{1} + \beta) + \varepsilon^{2}(H_{2} + H_{1}\beta + \tfrac12\beta\beta) + O(\varepsilon^{3})$, whose Einstein map has the coefficient $\mathcal{G}^{(1)}[H_{2} + H_{1}\beta + \tfrac12\beta\beta] + \mathcal{G}^{(2)}[H_{1} + \beta]$ of $\varepsilon^{2}$. By (1) it is also the coefficient of $\varepsilon^{2}$ of $U\,\mathcal{G}_{\varepsilon}(\tilde U a U)\,\tilde U$, which is $\mathcal{G}^{(1)}[H_{2}] + \mathcal{G}^{(2)}[H_{1}] + \mathcal{G}^{(1)}[H_{1}]\beta - \beta\,\mathcal{G}^{(1)}[H_{1}]$, since $U v \tilde U = v + \varepsilon\,B \cdot v = v - \varepsilon\,\beta(v)$ to first order. Subtracting gives the law. \emph{Displacement.} With $f(x) = x + \varepsilon\,\xi(x)$: $\underline{f} = \id + \varepsilon\,\underline{\xi}$, $\underline{f}^{-1} = \id - \varepsilon\,\underline{\xi} + \varepsilon^{2}\underline{\xi}\,\underline{\xi}$, and $\hb_{\varepsilon}(a; f(x)) = a + \varepsilon H_{1}(a) + \varepsilon^{2}((\xi \cdot \nabla)H_{1}(a) + H_{2}(a))$, all up to $O(\varepsilon^{3})$ and at $x$. So $\hb'_{\varepsilon} = \id + \varepsilon(H_{1} - \underline{\xi}) + \varepsilon^{2}(H_{2} + (\xi \cdot \nabla)H_{1} - \underline{\xi}H_{1} + \underline{\xi}\,\underline{\xi})$, with coefficient $\mathcal{G}^{(1)}[H_{2} + (\xi \cdot \nabla)H_{1} - \underline{\xi}H_{1} + \underline{\xi}\,\underline{\xi}] + \mathcal{G}^{(2)}[H_{1} - \underline{\xi}]$ of $\varepsilon^{2}$, and by (2) that coefficient is the one of $\mathcal{G}_{\varepsilon}(a; x + \varepsilon\xi)$, namely $\mathcal{G}^{(1)}[H_{2}] + \mathcal{G}^{(2)}[H_{1}] + (\xi \cdot \nabla)\mathcal{G}^{(1)}[H_{1}]$. Subtracting gives the law. \emph{Parametrisation.} Replacing $H_{2}$ by $H_{2} + q$ adds $\mathcal{G}^{(1)}[q]$ to the coefficient of $\varepsilon^{2}$ and nothing else; $(\id - \varepsilon H_{1} - \varepsilon^{2}H_{2})^{-1} = \id + \varepsilon H_{1} + \varepsilon^{2}(H_{2} + H_{1}H_{1})$ and $\exp(\varepsilon H_{1} + \varepsilon^{2}H_{2}) = \id + \varepsilon H_{1} + \varepsilon^{2}(H_{2} + \tfrac12 H_{1}H_{1})$ up to $O(\varepsilon^{3})$. On a first-order solution $\mathcal{G}^{(1)}[H_{1}] = 0$ and, by (i), the transformed $H_{1}$ is again a solution; the three laws reduce to the stated changes. \emph{Not invariant.} On a first-order solution replace $\xi$ by $s\,\xi$: $X_{s\xi} = s\,(-(\xi \cdot \nabla)H_{1} + \underline{\xi}H_{1}) - s^{2}\,\underline{\xi}\,\underline{\xi}$, so the change of $\mathcal{G}^{(2)}$ at a point is a polynomial in $s$ with the coefficient $-\mathcal{G}^{(1)}[\underline{\xi}\,\underline{\xi}]$ of $s^{2}$; where that coefficient is not zero the change is not zero for some $s$. $\mathcal{G}^{(1)}[\underline{\xi}\,\underline{\xi}]$ is a quadratic expression in the first, second and third derivatives of $\xi$ at the point with constant coefficients, the same at every point, and it is nonzero for the polynomial $\xi$ of the check at its point; translating that $\xi$ gives one for any point. The case $H_{1} = 0$ of the law is the displaced flat field: $H_{1} = -\underline{\xi}$ is a first-order solution with $\mathcal{G}^{(2)}[-\underline{\xi}] = -\mathcal{G}^{(1)}[\underline{\xi}\,\underline{\xi}]$. The check verifies both laws exactly off solutions on a polynomial $H_{1}$ that is not symmetric, with polynomial $\xi$ and $B$, and that leaving out any part of the quadratic field changes all sixteen components; on solutions it verifies that the change equals $\mathcal{G}^{(1)}[X]$ and is nonzero in all sixteen components; it computes the quadratic part of the displaced flat field from the gauge fields and finds $-\mathcal{G}^{(1)}[\underline{\xi}\,\underline{\xi}]$, nonzero in all sixteen components; and it verifies the parametrisation rule with a generic $q$.

\emph{(v).} For a field $x$ of symmetric maps, $\nabla \cdot (x(a)) = a \cdot d_{x}$, so $\nabla \cdot (\dot\nabla \wedge \dot x(a)) = \nabla^{2}x(a) - \nabla(a \cdot d_{x})$; for constant $a$ and a vector field $u$, $\nabla \cdot (a \wedge u) = (a \cdot \nabla)u - a\,(\nabla \cdot u)$. Hence $\nabla \cdot W_{a}[x] = \tfrac12\bigl(\nabla^{2}x(a) - \nabla(a \cdot d_{x}) - (a \cdot \nabla)d_{x} + a\,\nabla \cdot d_{x} + (a \cdot \nabla)\nabla\tau_{x} - a\,\nabla^{2}\tau_{x}\bigr) = \mathcal{E}[x](a)$. The divergence of the divergence of a bivector field vanishes. For the second identity, $Q(n) = 0$ gives $n \cdot \mathcal{E}[x](n) = \tfrac12\bigl(\nabla^{2}(n \cdot x(n)) + \partial_{n}^{2}\tau_{x} - 2\,\partial_{n}(n \cdot d_{x})\bigr)$. The reciprocals of $n, \bar n$ are $\tfrac12\bar n, \tfrac12 n$, so $\nabla = \tfrac12\bar n\,\partial_{n} + \tfrac12 n\,\partial_{\bar n} + \nabla_{\perp}$ with $\partial_{\bar n} = \bar n \cdot \nabla$, $\nabla^{2} = \partial_{n}\partial_{\bar n} + \nabla_{\perp}^{2}$, and $n \cdot d_{x} = \dot\nabla \cdot \dot x(n) = \tfrac12\,\partial_{n}(\bar n \cdot x(n)) + \tfrac12\,\partial_{\bar n}(n \cdot x(n)) + \nabla_{\perp} \cdot x(n)$. The two terms with $\partial_{n}\partial_{\bar n}$ cancel, leaving the stated expression, in which only derivatives along $N$ occur. \emph{Fluxes.} For a bivector field $W$ and the constant $n$, $n \cdot (\nabla \cdot W) = -\nabla \cdot (n \cdot W)$, and $u = n \cdot W$ has $n \cdot u = 0$, so $u = \alpha\,n + u_{\perp}$ with $u_{\perp}$ in the span of $e_{1}, e_{2}$, and $\nabla \cdot u = \partial_{n}\alpha + \nabla_{\perp} \cdot u_{\perp}$. The flux through $N$ of $\nabla \cdot W_{c}[x]$ is therefore $-\int_{N}(\partial_{\lambda}\alpha + \nabla_{\perp} \cdot u_{\perp})\,d\lambda\,d^{2}y$, which vanishes when $x$ vanishes near $N$ outside a compact part of $N$. By (iii), on a first-order solution $t(c)$ is a conserved vector field for each constant $c$; by (iv) a gauge change or a change of parametrisation adds $\tfrac{1}{8\pi G}\nabla \cdot W_{c}[-(X + X^{T})]$ to it; so $P_{c}$ is unchanged when $X$ has that support, which for $X_{\xi}$ and $X_{B}$ holds when $\xi$ or $B$ does. For $c = n$ the second identity gives the change of $P_{n}$ as $\tfrac{1}{8\pi G}$ times $\int d^{2}y\,\bigl[\tfrac12\,\partial_{n}(\tau_{x} - \bar n \cdot x(n)) - \nabla_{\perp} \cdot x(n)\bigr]$ between the two ends of the generators, plus the integral over $\lambda$ of a transverse divergence; the first bracket tends to zero where $\partial_{n}x$ does, and the remaining terms are transverse divergences. Both identities are proved by the computation above. In symbols (a momentum $p$ for each derivative) each side of each is linear in $x$ and quadratic in $p$, so it is fixed by its values on the ten unit symmetric maps and on the ten momenta $e_{i}$ and $e_{i} + e_{j}$, a quadratic expression in $p$ being fixed by its values on these; the check verifies both identities on those hundred pairs, which is a second proof, together with $p \cdot (p \cdot W) = 0$ and $n \cdot (n \cdot W) = 0$. The statements on fluxes are integrations and have no check.

\emph{(vi).} Let $H_{1}$ be symmetric with $H_{1}(n) = 0$ and take $H_{2} = r\,H_{1}H_{1}$, so that the coefficient of $\varepsilon^{2}$ of $\mathcal{G}_{\varepsilon}$ is $\mathcal{G}^{(2)}_{r}$. Then $\hb_{\varepsilon}^{-1} = \id + s$ with $s = -\varepsilon H_{1} + \varepsilon^{2}(1 - r)H_{1}H_{1} + O(\varepsilon^{3})$, a field of symmetric maps with $s(n) = 0$. By clause (iv) of the first cited theorem, which holds for such a field of any size, its clause (iii) applies with $K = n$ and $K \cdot D = \partial_{n}$ on scalars, and $\tfrac12\,\partial_{n}g(a, b) = \hb^{-1}(a) \cdot \mathsf{B}(\hb^{-1}(b))$. Since $g(a, b) = a \cdot (\id + s)^{2}(b)$, this is $\mathsf{B} = \tfrac12(\id + s)^{-1}\,\partial_{n}\bigl((\id + s)^{2}\bigr)\,(\id + s)^{-1}$, so $\tr\mathsf{B} = \tfrac12\tr\bigl((\id + s)^{-2}\partial_{n}(\id + s)^{2}\bigr) = \partial_{n}\tr s - \tfrac12\,\partial_{n}\tr(s^{2}) + O(s^{3})$ (with $m = 2s + s^{2}$: $\tfrac12\tr((\id - m)\partial_{n}m) = \tfrac12\partial_{n}\tr m - \tfrac14\partial_{n}\tr m^{2}$). By clause (iii)(c) there, $\theta = \tr\mathsf{B}$, $\tr\mathsf{B}^{2} = \tfrac12\theta^{2} + \ip{\sigma}{\sigma}$ and $d\theta/d\lambda = -\tr\mathsf{B}^{2} + K \cdot \mathcal{R}(K)$, exactly; and $n \cdot \mathcal{G}(n) = n \cdot \mathcal{R}(n)$ since $Q(n) = 0$. So $n \cdot \mathcal{G}_{\varepsilon}(n) = \partial_{n}\tr\mathsf{B} + \tr\mathsf{B}^{2}$. Now $\tr s = -\varepsilon\tr H_{1} + \varepsilon^{2}(1 - r)\tr(H_{1}H_{1})$, $\tr s^{2} = \varepsilon^{2}\tr(H_{1}H_{1})$ and $\mathsf{B} = -\varepsilon\,\partial_{n}H_{1} + O(\varepsilon^{2})$, so the coefficient of $\varepsilon^{2}$ is
\[
n \cdot \mathcal{G}^{(2)}_{r}(n) = \tr\bigl((\partial_{n}H_{1})^{2}\bigr) + \bigl(\tfrac12 - r\bigr)\,\partial_{n}^{2}\tr(H_{1}H_{1}) .
\]
With $h = -2H_{1}$: $\tr(H_{1}H_{1}) = \tfrac14\ip{h}{h}$, and $\tr((\partial_{n}H_{1})^{2}) = \tfrac14\ip{\partial_{n}h}{\partial_{n}h} = \tfrac14\ip{\partial_{n}h^{TT}}{\partial_{n}h^{TT}} + \tfrac18(\partial_{n}\tau)^{2} = \ip{\sigma}{\sigma} + \tfrac12\theta^{2}$, by the decomposition $\ip{h_{1}}{h_{2}} = \ip{h_{1}^{TT}}{h_{2}^{TT}} + \tfrac12\tau_{1}\tau_{2}$ of maps annihilating $n$ in the proof of clause (ii) of that theorem. This is the stated identity; no field equation was used. On a first-order solution, clause (ii) of that theorem gives $0 = n \cdot \mathcal{E}[h](n) = \tfrac12\,\partial_{n}^{2}\tau = \partial_{n}\theta$. Under (F), $\partial_{n}\tau$ tends to zero at the ends of a generator, so the constant $\theta$ is zero, and $\int\partial_{n}^{2}\ip{h}{h}\,d\lambda$ is the difference of $2\ip{h}{\partial_{n}h}$ between the ends, zero because $h$ is bounded and $\partial_{n}h$ tends to zero. The integral over $y$ gives $P_{n}$, and the right side is the expression $P$ of clause (vi) of that theorem; that it is the null energy, the generator of the null translations with the constant $8\pi G$, is that clause, under its hypotheses (L), (D), (V) and (N), which the equality of the integrals does not need. The last sentence of the clause is (iv) with (v), and (iv) gives a change of the form $\mathcal{G}^{(1)}[X]$ on solutions only. The check verifies the identity for $r = 0, 1, \tfrac12$ on two solutions in the gauge, each built from four null waves with $n \cdot p \neq 0$ and the displacement $\xi = \tfrac12\nu\,\bar n - \tfrac12(\bar n \cdot x)\nabla\nu$ with $\nu$ a polynomial of the transverse position, which preserves the gauge and gives $\theta \neq 0$; the three values differ, and two of the six are negative. It verifies the identity also on a polynomial field in the gauge that is not a solution.

\emph{(vii).} $\epsilon_{+}$ and $\epsilon_{\times}$ annihilate $n$ and $\bar n$ and have no trace, so $h(n) = 0$, $h(k) = 0$, $\tau = 0$; $h$ depends on $x$ through $k \cdot x$ only, so $\nabla^{2}h = Q(k)\,h'' = 0$, $d = h'(k) = 0$, and $\mathcal{E}[h] = 0$. By the Preliminaries the value of $\mathcal{G}^{(2)}_{r}$ at a point is a quadratic form in the values and the first and second derivatives of $H_{1}$ there, hence for this family a quadratic form in the six numbers $f_{+}, f_{+}', f_{+}'', f_{\times}, f_{\times}', f_{\times}''$ at the point, the same form at every point. So is the right side: $\ip{h'}{h'} = 2(f_{+}'^{2} + f_{\times}'^{2})$ and $\ip{h}{h}'' = 4(f_{+}'^{2} + f_{+}f_{+}'' + f_{\times}'^{2} + f_{\times}f_{\times}'')$. Two quadratic forms on six variables that agree on the six unit vectors and on the fifteen sums of two of them are equal. The check verifies that agreement for $r = \tfrac12$, all sixteen components on each of the 21 arguments; the cases $r = 0$ and $r = 1$ follow from (iv), $\mathcal{G}^{(2)}_{r} = \mathcal{G}^{(2)}_{1/2} + (r - \tfrac12)\,\mathcal{G}^{(1)}[H_{1}H_{1}]$ with $\mathcal{G}^{(1)}[H_{1}H_{1}] = \mathcal{E}[-\tfrac12 hh]$ and, for a field $x$ of the phase with $x(k) = 0$, $\mathcal{E}[x](a) = \tfrac12\,\tau_{x}''\,(a \cdot k)\,k$, here $-\tfrac14\ip{h}{h}''(a \cdot k)k$; the check confirms them on a generic argument. On $N$, $\partial_{n}h = (n \cdot k)h' = 2h'$, so $\sigma = h'$ and $n \cdot \tfrac{1}{32\pi G}\ip{h'}{h'}(n \cdot k)k = \tfrac{1}{8\pi G}\ip{\sigma}{\sigma}$, in agreement with (vi), as is the term $(\tfrac18 - \tfrac{r}{4})\,\partial_{n}^{2}\ip{h}{h} = (\tfrac12 - r)\ip{h}{h}''$.

\emph{(viii).} $(*)$ of (iii) was derived for any field $M_{\varepsilon}$ of maps with symmetric first-order part. With $M = T_{\varepsilon}$, (M2) says that its coefficients of $\varepsilon$ and of $\varepsilon^{2}$ vanish: $\nabla \cdot (T_{1}(c)) = 0$ and $\nabla \cdot (T_{2}(c)) = -F(c)$. By (iii) and (M1), $\nabla \cdot (\mathcal{G}^{(2)}[H_{1}](c))$ is $8\pi G$ times the same expression $F(c)$, so $\nabla \cdot (t(c)) = F(c)$ and the sum has no flat divergence. Conversely, with (M1), $\nabla \cdot (T_{1}(c)) = 0$ by (i) and $\nabla \cdot (t(c)) = F(c)$ by (iii), so $\nabla \cdot ((T_{2} + t)(c)) = 0$ is $\nabla \cdot (T_{2}(c)) = -F(c)$, the coefficient of $\varepsilon^{2}$ of (M2). The flat divergence of $\mathcal{G}^{(1)}[H_{2}]$ vanishes by (i), whence the condition on $\chi$; $\nabla \cdot ((T_{2} + \mu t)(c)) = (\mu - 1)F(c)$, and $\mathcal{G}^{(1)}[q]$ has no flat divergence. The check realises (M1) and (M2) with matter that is not built from the Einstein map: dust, $T(a) = \rho\,(a \cdot v)\,v$ with $\rho = \varepsilon\rho_{0} + \varepsilon^{2}\rho_{1}$ and $v = e_{0} + \varepsilon v_{1}$, for which $\operatorname{div}T(c) = (c \cdot v)\,D \cdot (\rho v) + \rho\,c \cdot (v \cdot D\,v)$, so that (M2) through second order is the geodesic equation at first order, $\partial_{t}v_{1} = -\Omega_{1}(e_{0}) \cdot e_{0}$, and the continuity equation at second. In the static field $h(a) = \varphi\,(2(a \cdot e_{0})e_{0} - a)$ one has $\mathcal{E}[h](a) = (\nabla^{2}\varphi)(a \cdot e_{0})e_{0}$, so (M1) holds with $8\pi G\,\rho_{0}$ the Laplacian of $\varphi$ on the rest space, and $\partial_{t}v_{1} = \tfrac12\nabla\varphi$, Newton's law for the potential $\tfrac12\varphi$. With polynomial $\varphi$, $v_{1}$ and $\rho_{1}$ are polynomials; the check finds the flat divergence of $T_{2}$ nonzero and that of $T_{2} + t$ zero, exactly, in the symmetric static gauge and after a polynomial displacement and rotation gauge change (under which $T_{2}$ changes by $(\xi \cdot \nabla)T_{1} + T_{1}\beta - \beta T_{1}$, by the covariance of (iv)), where all four components of $F$ are nonzero; the force being nonzero, no other constant in front of the quadratic part balances.

\emph{Reading, the instance.} Let $\ell = \tfrac12\bar n$ ($Q(\ell) = 0$, $n \cdot \ell = 1$), $\lambda = \ell \cdot x$, $y$ the part of $x$ in the span of $e_{1}, e_{2}$, $\mathsf{w}_{\lambda}$ a family of symmetric maps of that span with $\mathsf{w}(n) = \mathsf{w}(\ell) = 0$, $W = -y \cdot \mathsf{w}_{\lambda}(y)$ and $\hb_{\varepsilon}(a) = a + \tfrac12\varepsilon W (a \cdot \ell)\,\ell$. Since $Q(\ell) = 0$, $g_{b} = b - \tfrac12\varepsilon W (b \cdot \ell)\,\ell$ exactly, and the square of the map $a \mapsto \tfrac12 W (a \cdot \ell)\ell$ is zero. Put $\Omega(a) = \varepsilon\,(a \cdot \ell)\,\Psi$ with $\Psi = \ell \wedge \mathsf{w}(y)$. Then $\Psi \cdot \ell = 0$ and $\Psi \cdot b = (b \cdot \mathsf{w}(y))\ell - (b \cdot \ell)\mathsf{w}(y)$, and $\nabla W = -2\mathsf{w}(y) - (y \cdot \mathsf{w}'(y))\ell$, so $D_{a}g_{b} = \varepsilon\bigl((a \cdot \mathsf{w}(y))(b \cdot \ell) + (a \cdot \ell)(b \cdot \mathsf{w}(y))\bigr)\ell + \tfrac12\varepsilon(y \cdot \mathsf{w}'(y))(a \cdot \ell)(b \cdot \ell)\ell - \varepsilon(a \cdot \ell)(b \cdot \ell)\mathsf{w}(y)$, symmetric in $a, b$: (T2) holds exactly, and this is the $\Omega$ of the field. $\Omega(a) \times \Omega(b)$ is a multiple of $\Psi \times \Psi = 0$, and $(a \cdot \nabla)\Omega(b) = \varepsilon(b \cdot \ell)\,\ell \wedge \bigl((a \cdot \ell)\mathsf{w}'(y) + \mathsf{w}(a)\bigr)$, so $R(a \wedge b) = \varepsilon\,\ell \wedge \mathsf{w}((a \wedge b) \cdot \ell)$; and $\hb(a \wedge b) \cdot \ell = (a \wedge b) \cdot \ell$ because $\hb(a) \cdot \ell = a \cdot \ell$. Hence $\mathcal{R}(B) = -\varepsilon\,\ell \wedge \mathsf{w}(B \cdot \ell)$ and, contracting, $\mathcal{G}_{\varepsilon}(a) = \varepsilon\,(\tr\mathsf{w})\,(a \cdot \ell)\,\ell$, with no term of higher order: the quadratic part vanishes at every point, for every $\mathsf{w}$, and for every $r$ since the square of $H_{1}$ is zero; and $H_{1}$ is a first-order solution exactly when $\tr\mathsf{w} = 0$. \emph{The displacement.} Let $F_{\lambda}$ be symmetric maps of the span of $e_{1}, e_{2}$ with $F(n) = F(\ell) = 0$ and $F'' = 2\mathsf{w}$, and $\xi = -\tfrac12 F(y) + \tfrac14(y \cdot F'(y))\,\ell$. With $\nabla\lambda = \ell$ and $(a \cdot \nabla)y$ the transverse part of $a$: the map $a \mapsto (a \cdot \nabla)F(y) + \nabla(a \cdot F(y))$ is $2F(a) + (a \cdot \ell)F'(y) + (a \cdot F'(y))\ell$, and for $u = (y \cdot F'(y))\ell$ the map $a \mapsto (a \cdot \nabla)u + \nabla(a \cdot u)$ is $2(y \cdot F''(y))(a \cdot \ell)\ell + 2(a \cdot F'(y))\ell + 2(a \cdot \ell)F'(y)$. So $\underline{\xi} + \underline{\xi}^{T} = -F + \tfrac12(y \cdot F''(y))\,(a \mapsto (a \cdot \ell)\ell)$, and $-\tfrac12 F - \tfrac12(\underline{\xi} + \underline{\xi}^{T})$ is the map $a \mapsto \tfrac12 W (a \cdot \ell)\ell$: the Kerr-Schild $H_{1}$ is $-\tfrac12 F - \underline{\xi} + \beta_{B}$ with $\beta_{B} = \tfrac12(\underline{\xi} - \underline{\xi}^{T})$. By (iv), applied to the displacement and then to the rotation, the two quadratic parts differ by $\mathcal{G}^{(1)}[X]$ with $X = X_{\xi}[-\tfrac12 F] + X_{B}[-\tfrac12 F - \underline{\xi}]$. For trace-free $F$ the wave $h = F$ is the one of (vii) with phase $k \cdot x = 2\lambda$, whose quadratic part is not zero. The check builds both fields with polynomial profiles and verifies, exactly: the torsion equation, the vanishing square, the vanishing of the whole jet of the first-order part and of the quadratic part of the Kerr-Schild field; with a trace, the first-order part $(\tr\mathsf{w})(a \cdot \ell)\ell$ and again no quadratic part; the relation between the two fields; that the quadratic part of the transverse trace-free wave is nonzero and equal to minus $\mathcal{G}^{(1)}[X]$; and that $\xi$ has nonzero second derivatives across the plane. Under (F) for $F$, (vi) gives, on each generator, $\int n \cdot t(n)\,d\lambda = \tfrac{1}{8\pi G}\int\ip{\sigma}{\sigma}\,d\lambda$ with $\sigma = \tfrac12 F'$ for the first field, and $t = 0$ gives zero for the second; $P_{n}$ over the whole plane diverges for a plane wave.
